\documentclass[11pt,a4paper]{article}

% ===== PACKAGES =====
\usepackage[utf8]{inputenc}
\usepackage[T1]{fontenc}
\usepackage{geometry}
\geometry{margin=2.5cm}
\usepackage{times}
\usepackage{graphicx}
\usepackage{booktabs}
\usepackage{array}
\usepackage{multirow}
\usepackage{enumitem}
\usepackage{xcolor}
\usepackage{hyperref}
\usepackage{tabularx}
\usepackage{caption}
\usepackage{float}
\usepackage{amsmath}
\usepackage{tcolorbox}
\usepackage{fancyhdr}

% ===== COLORS =====
\definecolor{cmuRed}{RGB}{196,18,48}
\definecolor{darkBlue}{RGB}{25,60,120}
\definecolor{noteBlue}{RGB}{232,244,253}
\definecolor{noteBorder}{RGB}{66,133,244}
\definecolor{warnOrange}{RGB}{255,243,224}
\definecolor{warnBorder}{RGB}{230,126,34}
\definecolor{greenBox}{RGB}{232,250,237}
\definecolor{greenBorder}{RGB}{39,174,96}

% ===== CUSTOM BOXES =====
\newtcolorbox{notebox}{
  colback=noteBlue, colframe=noteBorder, 
  boxrule=0.5pt, left=6pt, right=6pt, top=4pt, bottom=4pt,
  fontupper=\small\itshape
}

\newtcolorbox{warnbox}{
  colback=warnOrange, colframe=warnBorder,
  boxrule=0.5pt, left=6pt, right=6pt, top=4pt, bottom=4pt,
  fontupper=\small\itshape
}

\newtcolorbox{greenbox}{
  colback=greenBox, colframe=greenBorder,
  boxrule=0.5pt, left=6pt, right=6pt, top=4pt, bottom=4pt,
  fontupper=\small\itshape
}

% ===== HEADER/FOOTER =====
\pagestyle{fancy}
\fancyhf{}
\fancyhead[C]{\small\textit{Addressing Gaps in LLM-Based Learning Objective Generation}}
\fancyfoot[C]{\thepage}
\renewcommand{\headrulewidth}{0.4pt}

% ===== TITLE =====
\title{
\vspace{-1cm}
\textbf{Beyond Single-Shot Prompting:}\\[0.3cm]
\large Improving the Robustness, Generalizability, and Evaluation of LLM-Generated Learning Objectives\\[0.5cm]
\normalsize A Gap Analysis \& Extension Proposal
}

\author{
ProfsAgent Research Notes\\
\textit{Building upon: Sridhar et al. (2023) \cite{sridhar2023}}
}

\date{August 2026}

% ===== DOCUMENT =====
\begin{document}

\maketitle
\thispagestyle{fancy}

% ---------- ABSTRACT ----------
\begin{abstract}
Sridhar et al.\ (2023) demonstrated that GPT-4 can generate sensible, well-structured Learning Objectives (LOs) aligned with Bloom's Taxonomy for a university-level AI course. However, their work relies on a single model (GPT-4), a single temperature setting (0.7), a single course domain (AI Practitioner at CMU), and a narrow evaluation panel of 3 CS graduate students achieving only fair inter-rater agreement ($\kappa = 0.31$). We identify \textbf{12 key gaps} in their methodology, evaluation, and generalizability, and propose concrete experiments to address them---including temperature-diversity analysis, multi-model benchmarking across 6 LLMs, region-specific Bloom's classifiers for globally applicable validation, and a multi-perspective annotator panel comprising professors and students of varying skill levels. These extensions aim to improve the quality, generalizability, and reliability of automated LO generation for real-world educational deployment.
\end{abstract}

\noindent\rule{\textwidth}{0.4pt}
\tableofcontents
\noindent\rule{\textwidth}{0.4pt}
\vspace{0.3cm}

% ======================================================================
\section{Introduction}
% ======================================================================

Learning Objectives (LOs) are the blueprints against which all course content, instructional strategies, and assessments are designed \cite{bloom1956, krathwohl2002}. Despite their importance, authoring high-quality LOs remains a time-consuming and expertise-intensive task that many instructors skip or perform inadequately. Sridhar et al.\ \cite{sridhar2023} explored using GPT-4 to automate LO generation, demonstrating promising results: 127 LOs generated for an AI Practitioner course were found to be largely sensible, properly expressed with action verbs, and aligned with Bloom's Taxonomy.

However, upon close analysis, we identify \textbf{12 significant gaps} in the study's methodology, evaluation design, and generalizability. This paper systematically documents each gap and proposes concrete experiments and extensions to address them, with the goal of building a more robust, globally applicable framework for LLM-assisted LO generation.

% ======================================================================
\section{Gap 1: No Temperature / Hyperparameter Variation}
% ======================================================================

\subsection{Problem}
The paper uses \texttt{temperature=0.7} (the default) with no justification or experimentation. Temperature directly controls the creativity--coherence trade-off---a critical factor for LO diversity and quality. No other hyperparameters (\texttt{top\_p}, penalties) were varied either.

\subsection{Proposed Experiment}

\begin{table}[H]
\centering
\caption{Proposed temperature sweep experimental design.}
\label{tab:temp}
\begin{tabular}{@{}clp{5.5cm}@{}}
\toprule
\textbf{Temperature} & \textbf{Expected Behavior} & \textbf{What to Measure} \\
\midrule
0.2 & Highly deterministic, repetitive & Quality $\uparrow$, Diversity $\downarrow$ \\
0.5 & Balanced, conservative & Baseline comparison \\
0.7 & Default (Paper 1's choice) & Replicate original results \\
0.9 & Creative, more varied & Diversity $\uparrow$, coherence risk \\
1.2 & Highly creative, noisy & Risk of nonsensical LOs \\
\bottomrule
\end{tabular}
\end{table}

\subsection{Metrics}
\begin{itemize}[nosep]
    \item \textbf{Quality:} Sensibility score (Likert 1--5), Bloom's alignment accuracy
    \item \textbf{Diversity:} Unique action verbs per module, pairwise BLEU/BERTScore between LOs (lower = more diverse)
    \item \textbf{Best-of-N:} Generate $N=5$ sets at each temperature, select best set via a quality scorer
\end{itemize}

\begin{notebox}
\textbf{Expected Contribution:} Identify the optimal temperature range and demonstrate that temperature tuning alone can measurably improve LO quality and diversity.
\end{notebox}

% ======================================================================
\section{Gap 2: Single Model Benchmarking (GPT-4 Only)}
% ======================================================================

\subsection{Problem}
Only GPT-4 was tested. No comparison with any other model---not even GPT-3.5. The paper was published in June 2023; by 2026, the LLM landscape includes GPT-4o, Claude, Gemini, Llama~3, Mistral, and many more capable models across both commercial and open-source categories.

\subsection{Proposed Multi-Model Benchmark}

\begin{table}[H]
\centering
\caption{Proposed models for multi-model LO generation benchmark.}
\label{tab:models}
\begin{tabular}{@{}llcp{5cm}@{}}
\toprule
\textbf{Model} & \textbf{Type} & \textbf{Cost} & \textbf{Why Include} \\
\midrule
GPT-4o & Commercial & \$\$\$ & Current OpenAI flagship \\
GPT-4o-mini & Commercial & \$ & Cost-effective alternative \\
Claude Sonnet 4 & Commercial & \$\$ & Strong reasoning, different training data \\
Gemini 2.5 Pro & Commercial & \$\$ & Google's flagship model \\
Llama 3.1 70B & Open-source & Free & Privacy---no data leaves institution \\
Mistral Large & Open-source & Free & European open-source alternative \\
\bottomrule
\end{tabular}
\end{table}

\subsection{Evaluation Protocol}
Run the \textit{exact same prompt} (system + user) from Paper~1 across all models. Compare on RQ1 (sensibility), RQ2 (action verbs), RQ3 (Bloom's alignment), plus cost per LO, latency, and specificity of outputs.

\begin{notebox}
\textbf{Expected Contribution:} First multi-model benchmark for educational LO generation. May reveal that cheaper or open-source models perform comparably to GPT-4, enabling cost-effective deployment.
\end{notebox}

% ======================================================================
\section{Gap 3: Region-Specific BERT Classifier}
% ======================================================================

\subsection{Problem}
The BERT classifier used for automated Bloom's level classification \cite{li2022} was trained on 21,380 LOs from 5,558 courses---almost certainly from US/English-speaking universities. Educational conventions vary dramatically by country and accreditation body. A classifier trained on US-style LOs will likely misclassify LOs written in Indian, European, or other regional educational conventions.

\subsection{Regional Differences in LO Conventions}

\begin{table}[H]
\centering
\caption{LO conventions vary significantly across educational systems.}
\label{tab:regions}
\begin{tabular}{@{}lp{3.2cm}p{3.2cm}p{3.5cm}@{}}
\toprule
\textbf{Aspect} & \textbf{US Convention} & \textbf{Indian Convention} & \textbf{European Convention} \\
\midrule
LO Format & Student-centered: ``Students will be able to\ldots'' & CO-based: ``CO1: Design algorithms\ldots'' & Competency-based: ``Graduate can demonstrate\ldots'' \\
\addlinespace
Taxonomy & Bloom's (6 levels) & Bloom's + SOLO taxonomy & Dublin Descriptors / EQF levels \\
\addlinespace
Accreditation & ABET & NBA / NAAC & EUR-ACE / Bologna Process \\
\addlinespace
Style & Informal, specific & Formal, sometimes broader & Outcome-oriented, transferable skills \\
\bottomrule
\end{tabular}
\end{table}

\subsection{Proposed Solution}
\begin{enumerate}[nosep]
    \item Collect LO datasets from universities across \textbf{India} (NBA/NAAC), \textbf{Europe} (Bologna/EUR-ACE), and other regions.
    \item Fine-tune \textbf{separate BERT classifiers} or one \textbf{multi-region classifier} with region labels.
    \item Allow users to select their country/accreditation body, routing to the appropriate classifier.
    \item Evaluate cross-region classification accuracy vs.\ single-region baselines.
\end{enumerate}

\begin{notebox}
\textbf{Expected Contribution:} First region-aware Bloom's taxonomy classifier, enabling globally applicable LO validation beyond US-centric educational norms.
\end{notebox}

% ======================================================================
\section{Gap 4: Narrow, Homogeneous Annotator Panel}
% ======================================================================

\subsection{Problem}
The paper used 3 CS graduate students---all from the same institution (CMU), same expertise level, same perspective. They achieved only $\kappa = 0.31$ (``fair'' agreement). This low agreement partly stems from having a homogeneous panel with no clearly defined evaluation criteria per role.

\subsection{Proposed Multi-Perspective Annotator Panel}

\begin{table}[H]
\centering
\caption{Proposed diverse annotator panel for LO evaluation.}
\label{tab:annotators}
\begin{tabular}{@{}lcll@{}}
\toprule
\textbf{Role} & \textbf{Count} & \textbf{Evaluates} & \textbf{Unique Perspective} \\
\midrule
Professor / Expert & 1 & Bloom's accuracy, pedagogy, accreditation & ``Is this what I'd teach?'' \\
Strong Students (Top 20\%) & 3 & Challenge level, advanced topics & ``Is this too easy?'' \\
Average Students (Mid 40\%) & 3 & Clarity, accessibility, expectations & ``Do I understand this?'' \\
\midrule
\textbf{Total} & \textbf{7} & & \textbf{3 distinct perspectives} \\
\bottomrule
\end{tabular}
\end{table}

\subsection{Why This Improves $\kappa$}
\begin{itemize}[nosep]
    \item The paper's $\kappa = 0.31$ partly results from all 3 annotators sharing the same background---their disagreements were noisy rather than informative.
    \item A structured panel with \textbf{defined evaluation criteria per role} reduces ambiguity.
    \item The professor provides an \textbf{expert anchor}; students validate from the \textbf{learner's perspective}.
    \item Reporting $\kappa$ \textbf{per role-group} reveals where LOs fail for different stakeholders.
\end{itemize}

\begin{warnbox}
\textbf{Key Insight:} Average students are arguably the \textbf{most important} validators. If an LO isn't clear to them, it fails its primary purpose---guiding the learner. The paper's annotator panel completely lacks this perspective.
\end{warnbox}

% ======================================================================
\section{Gap 5: No Baseline Comparison with Human-Written LOs}
% ======================================================================

\subsection{Problem}
The paper evaluates GPT-4's LOs in isolation. There is no comparison against LOs written by \textbf{actual professors} for the same modules. Without a baseline, we cannot answer the fundamental question: \textit{``Are GPT-4's LOs better, worse, or comparable to human-written ones?''}

\subsection{Proposed Experiment}
\begin{enumerate}[nosep]
    \item Collect professor-written LOs for the same AI Practitioner course modules.
    \item \textbf{Blind-evaluate} both sets (GPT-4 vs.\ human) using an identical rubric---evaluators should not know which set was machine-generated.
    \item Report head-to-head comparison on sensibility, Bloom's alignment, specificity, and measurability.
\end{enumerate}

\begin{notebox}
\textbf{Expected Contribution:} First direct GPT vs.\ human LO quality comparison---the most fundamental missing benchmark in the original study.
\end{notebox}

% ======================================================================
\section{Gap 6: No Iterative Refinement / Self-Critique}
% ======================================================================

\subsection{Problem}
Paper~1 uses single-shot generation---one API call, accept whatever GPT-4 outputs. No self-review, no refinement, no chain-of-thought reasoning. Modern LLM best practices demonstrate that self-critique significantly improves output quality.

\subsection{Proposed Approach}
\begin{enumerate}[nosep]
    \item \textbf{Generate:} Produce initial LOs using the same prompt as Paper~1.
    \item \textbf{Self-Critique:} Prompt the LLM to review its own LOs: \textit{``Review these LOs against the quality criteria. Flag issues with specificity, multiple action verbs, or Bloom's misalignment.''}
    \item \textbf{Refine:} Prompt the LLM to fix flagged issues and output improved LOs.
    \item \textbf{Validate:} Run the same evaluation (RQ1--RQ3) on both single-shot and refined LOs.
\end{enumerate}

\begin{notebox}
\textbf{Expected Contribution:} Demonstrate that a simple 2-step self-critique loop measurably improves LO quality over single-shot generation, with minimal additional cost.
\end{notebox}

% ======================================================================
\section{Gap 7: No Measurability Evaluation}
% ======================================================================

\subsection{Problem}
The paper explicitly states that LOs should be \textit{``measurable''}---but never evaluates whether the generated LOs actually are measurable. The authors themselves acknowledge this as future work.

\subsection{Proposed Rubric}

\begin{table}[H]
\centering
\caption{Proposed measurability evaluation rubric for LOs.}
\label{tab:measurability}
\begin{tabular}{@{}lp{4.5cm}p{4.5cm}@{}}
\toprule
\textbf{Criterion} & \textbf{Score 1 (Poor)} & \textbf{Score 5 (Excellent)} \\
\midrule
Observable Behavior & Cannot observe student performing this & Clear, visible student action \\
Assessment-Mappable & Cannot design a test for this & Direct test question is possible \\
Success Criteria & No standard for ``how well'' & Specific, quantifiable threshold \\
\bottomrule
\end{tabular}
\end{table}

\textbf{Examples:}
\begin{itemize}[nosep]
    \item \textcolor{red}{\textbf{Not measurable:}} \textit{``Understand neural networks''} --- How do you test ``understanding''?
    \item \textcolor{greenBorder}{\textbf{Measurable:}} \textit{``Explain the forward propagation process in a 3-layer neural network, identifying the role of weights, biases, and activation functions''} --- Can be directly tested.
\end{itemize}

\begin{notebox}
\textbf{Expected Contribution:} First measurability evaluation of LLM-generated LOs, directly addressing the paper's stated future work.
\end{notebox}

% ======================================================================
\section{Gap 8: No Prompt Component Ablation}
% ======================================================================

\subsection{Problem}
The system prompt contains multiple components: role definition, LO structure guidelines, 30 few-shot examples, and 4 quality criteria. Which components actually contribute to LO quality? The paper never investigates this.

\subsection{Proposed Ablation Study}

\begin{table}[H]
\centering
\caption{Proposed prompt ablation conditions.}
\label{tab:ablation}
\begin{tabular}{@{}llp{5cm}@{}}
\toprule
\textbf{Condition} & \textbf{What's Removed} & \textbf{Research Question} \\
\midrule
Full prompt & Nothing & Baseline (replicate Paper~1) \\
No examples & All 30 example LOs & Do few-shot examples matter? \\
No role & ``Curricular expert'' role & Does role-playing improve quality? \\
No criteria & 4 quality criteria & Do explicit rules matter? \\
Minimal & Only ``Generate LOs for [module]'' & Total prompt engineering contribution? \\
\bottomrule
\end{tabular}
\end{table}

\begin{notebox}
\textbf{Expected Contribution:} Quantify the contribution of each prompt component to LO quality---guiding optimal prompt design for educational applications.
\end{notebox}

% ======================================================================
\section{Gap 9: No RAG / Grounding on Actual Course Content}
% ======================================================================

\subsection{Problem}
LOs are generated from just a module name (e.g., ``Generative Models''). The model has no access to the actual lecture content, textbook chapters, or syllabus---leading to generic, sometimes vague LOs. The paper itself identifies the specificity issue (e.g., ``Python libraries'' instead of ``scikit-learn'').

\subsection{Proposed Approach}
Use \textbf{Retrieval-Augmented Generation (RAG)} to ground LO generation in actual course materials:
\begin{enumerate}[nosep]
    \item Retrieve relevant textbook sections, lecture notes, or prior syllabi via vector similarity search.
    \item Include retrieved content as context alongside the module name in the prompt.
    \item Compare RAG-grounded LOs vs.\ module-name-only LOs on specificity and relevance.
\end{enumerate}

\begin{notebox}
\textbf{Expected Contribution:} Demonstrate that grounding LOs in actual course content produces more specific, actionable, and context-appropriate objectives.
\end{notebox}

% ======================================================================
\section{Gap 10: No Cross-Module Consistency Check}
% ======================================================================

\subsection{Problem}
Each module's LOs are generated independently via separate API calls. There is no verification for:
\begin{itemize}[nosep]
    \item \textbf{Overlap:} Two modules generating redundant LOs covering the same skill.
    \item \textbf{Gaps:} Important Bloom's levels not covered across the entire course.
    \item \textbf{Prerequisite violations:} Higher-level LOs appearing before foundational ones in the course sequence.
\end{itemize}

\subsection{Proposed Solution}
After generating all module-level LOs, run a \textbf{course-level consistency analysis} that produces a coverage matrix (Modules $\times$ Bloom's Levels), flags duplicate LOs across modules, and verifies prerequisite ordering across the course timeline.

\begin{notebox}
\textbf{Expected Contribution:} Course-level LO coherence validation---moving beyond individual LO quality to holistic curriculum design.
\end{notebox}

% ======================================================================
\section{Gap 11: No Structured Output Format}
% ======================================================================

\subsection{Problem}
LOs are generated as free text, making automated downstream processing (Bloom's classification, alignment checking, syllabus mapping) difficult and error-prone.

\subsection{Proposed Solution}
Force structured JSON output with explicit fields for each LO:

\begin{verbatim}
{
  "action_verb": "Implement",
  "behavior": "a classification pipeline",
  "conditions": "using scikit-learn on a labeled dataset",
  "degree": "achieving >= 85% accuracy on test data",
  "blooms_level": "Apply",
  "module_type": "project"
}
\end{verbatim}

This enables automatic validation, Bloom's classification, downstream syllabus/assessment integration, and consistent quality checking across the entire course.

\begin{notebox}
\textbf{Expected Contribution:} Enable automated validation and downstream pipeline integration for scalable LO processing.
\end{notebox}

% ======================================================================
\section{Gap 12: No Downstream Evaluation of Learning Outcomes}
% ======================================================================

\subsection{Problem}
The paper evaluates LOs in isolation. The ultimate question remains unanswered: \textit{``Do courses designed around LLM-generated LOs actually produce better learning outcomes?''} This is the gold-standard evaluation that would truly validate the approach.

\subsection{Proposed Approach (Long-term)}
\begin{enumerate}[nosep]
    \item Deploy LLM-generated LOs in a real course for one semester.
    \item Compare student performance, satisfaction, and learning gains against a prior semester using human-written LOs.
    \item Run pre/post assessments aligned to the LOs to measure actual learning outcomes.
\end{enumerate}

\begin{warnbox}
\textbf{Note:} This is a longer-term study requiring IRB approval and institutional support, but it represents the ultimate validation of the entire approach. It should be framed as future work.
\end{warnbox}

% ======================================================================
\section{Summary and Prioritization}
% ======================================================================

Table~\ref{tab:summary} provides an overview of all 12 identified gaps along with their estimated effort and potential impact.

\begin{table}[H]
\centering
\caption{Summary of all 12 identified gaps with effort and impact assessment.}
\label{tab:summary}
\begin{tabular}{@{}clcc@{}}
\toprule
\textbf{\#} & \textbf{Gap Description} & \textbf{Effort} & \textbf{Impact} \\
\midrule
1 & Temperature / hyperparameter variation & Low & High \\
2 & Multi-model benchmarking (6 LLMs) & Low & High \\
3 & Region-specific BERT classifier & Medium & Very High \\
4 & Diverse annotator panel (7 people) & Medium & Very High \\
5 & Baseline comparison with human LOs & Low & Very High \\
6 & Iterative self-critique loop & Low & High \\
\midrule
7 & Measurability evaluation rubric & Medium & High \\
8 & Prompt component ablation study & Low & Medium \\
9 & RAG-grounded LO generation & Medium & High \\
10 & Cross-module consistency checking & Medium & Medium \\
11 & Structured output format (JSON) & Low & Medium \\
\midrule
12 & Downstream learning outcome evaluation & High & Very High \\
\bottomrule
\end{tabular}
\end{table}

\subsection{Recommended Prioritization}

\begin{greenbox}
\textbf{HIGH PRIORITY --- Include in Paper (Gaps 1--6):} These are experimentally feasible, require low-to-medium effort, and directly address the most critical weaknesses in the original study. They form the core contribution of this extension work.
\end{greenbox}

\begin{notebox}
\textbf{MEDIUM PRIORITY --- Discuss in Paper (Gaps 7--11):} These are important extensions that require more infrastructure (datasets, retrieval systems, rubric development) but should be discussed and partially explored.
\end{notebox}

\begin{warnbox}
\textbf{FUTURE WORK --- Mention (Gap 12):} The downstream evaluation requires a longitudinal study with real students across semesters. It should be clearly articulated as the natural next step for validating the entire approach.
\end{warnbox}

% ======================================================================
% REFERENCES
% ======================================================================
\bibliographystyle{unsrt}

\begin{thebibliography}{10}

\bibitem{sridhar2023}
P.~Sridhar, A.~Doyle, A.~Agarwal, C.~Bogart, J.~Savelka, and M.~Sakr.
\newblock Harnessing {LLMs} in curricular design: Using {GPT-4} to support authoring of learning objectives.
\newblock In \textit{Workshop on Empowering Education with LLMs, AIED}, 2023.
\newblock arXiv:2306.17459.

\bibitem{li2022}
Y.~Li, M.~Rakovic, B.~X.~Poh, D.~Gasevic, and G.~Chen.
\newblock Automatic classification of learning objectives based on {Bloom's} taxonomy.
\newblock In \textit{Proceedings of EDM}, pages 530--537, 2022.

\bibitem{bloom1956}
B.~S. Bloom, M.~D. Engelhart, E.~Furst, W.~H. Hill, and D.~R. Krathwohl.
\newblock \textit{Taxonomy of Educational Objectives: Handbook {I} --- Cognitive Domain}.
\newblock David McKay, New York, 1956.

\bibitem{krathwohl2002}
D.~R. Krathwohl.
\newblock A revision of {Bloom's} taxonomy: An overview.
\newblock \textit{Theory into Practice}, 41:212--218, 2002.

\bibitem{tran2018}
K.~N. Tran, J.~H. Lau, D.~Contractor, U.~Gupta, B.~Sengupta, C.~J. Butler, and M.~Mohania.
\newblock Document chunking and learning objective generation for instruction design.
\newblock In \textit{Proceedings of EDM}, 2018.

\bibitem{yao2026}
H.~Yao, W.~Xu, J.~Turnau, N.~Kellam, and H.~Wei.
\newblock Instructional agents: Reducing teaching faculty workload through multi-agent instructional design.
\newblock Arizona State University, 2026.
\newblock arXiv:2508.19611.

\bibitem{adams2015}
N.~E. Adams.
\newblock {Bloom's} taxonomy of cognitive learning objectives.
\newblock \textit{Journal of the Medical Library Association}, 103(3):152--153, 2015.

\end{thebibliography}

\end{document}
